
\documentclass[11pt]{article}
\usepackage[margin=1in]{geometry}

% Core packages
\usepackage{amsmath,amssymb}
\usepackage{tikz-cd}
\usepackage{multicol}
\usepackage{tabularx}

% Paragraphs
\setlength{\parindent}{0pt}
\setlength{\parskip}{1\baselineskip}


\begin{document}
\begin{titlepage}
\centering
\vspace*{1cm}
{\Large\textbf{Universidad de Antioquia}\par}
{\large Facultad de Ciencias Exactas y Naturales\par Instituto de Matemáticas\par}
\vspace{2cm}
{\LARGE\bfseries Aplicación de la estadística descriptiva al conjunto de datos \texttt{survey}\par}
\vspace{1.5cm}
{\large Nicol Dayana Vasquez Zapata\par}
{\normalsize Estudiante del Pregrado de Estadística\par}
\vspace{0.7cm}
{\large Daniel Soto Villada\par}
{\normalsize Estudiante del Pregrado de Astronomía\par}
\vfill
{\normalsize Curso: Fundamentos de Estadística I\par Docente: Luz Estela del Socorro Sánchez Herrera\par}
\vspace{0.75cm}
Semestre 2026-2\par
Fecha: \today\par
\end{titlepage}

\section{Introducción}

\section{Materiales y métodos}
\subsection{Descripción de la base de datos}

El conjunto de datos \texttt{survey}, incluido en el paquete \texttt{MASS} de R, contiene las respuestas de una encuesta realizada a \textbf{237 estudiantes de primer año de estadística} de la Universidad de Adelaida (Australia). Se utiliza con frecuencia en la enseñanza de la estadística para ilustrar el análisis exploratorio de datos, las pruebas de hipótesis —como la prueba de independencia $\chi^2$ y la prueba $t$ de Student— y los modelos de regresión.

\subsection{}

\begin{center}
\renewcommand{\arraystretch}{1.2}
\begin{tabularx}{\textwidth}{|l|X|X|}
\hline
\textbf{Variable} & \textbf{Descripción} & \textbf{Tipo de variable estadística} \\
\hline
\texttt{Sex} & Sexo reportado por el estudiante (\texttt{Female}, \texttt{Male}). & \textbf{Cualitativa nominal} \\
\hline
\texttt{Wr.Hnd} & Amplitud o envergadura de la mano con la que escribe (\textit{writing hand span}), en centímetros; distancia del pulgar al meñique extendido. & \textbf{Cuantitativa continua} \\
\hline
\texttt{NW.Hnd} & Amplitud o envergadura de la mano con la que no escribe (\textit{non-writing hand span}), en centímetros. & \textbf{Cuantitativa continua} \\
\hline
\texttt{W.Hnd} & Mano dominante para escribir (\texttt{Left} para zurdo, \texttt{Right} para diestro). & \textbf{Cualitativa nominal} \\
\hline
\texttt{Pulse} & Frecuencia cardíaca en reposo, en latidos por minuto (bpm). & \textbf{Cuantitativa discreta} (usualmente tratada como continua en modelos de regresión). \\
\hline
\texttt{Height} & Estatura del estudiante, expresada en centímetros. & \textbf{Cuantitativa continua} \\
\hline
\end{tabularx}
\end{center}

\subsection{Métodos estadísticos}
definición y fórmulas de las medidas de
resumen usadas (de tendencia central, dispersión, asimetría y apuntalamiento) y de sus valores de interpretación (por ejemplo, si $K> 0$ entonces..., si $K < 0$, entonces..., y si $K = 0$ entonces...).

\subsection{Descripción de valores perdidos}
Descripción de los valores perdidos: (Si los hay) se debe contar cuáles valores perdidos se encontraron y cómo se ven en la base de datos. Ejemplo: NC=No contestó, SD=Sin dato, NS=No sabe, etc.

\subsection{Descripción de los análisis}
Que es lo que se quiere observar en la base de datos

\section{Resultados}
Tablas, graficos y analisis correspondientes
\subsection{Análisis Univariado}
Para las variables cuantitativas calcular las medidas de resumen. Para variables cualitativas, hacer tablas de frecuencia (con n y \%) y diagrama de barras.
\subsection{Otros análisis cuantitativos}
Escoja dos o tres variable cuantitativa y muestre
claramente sus cuartiles y deciles, diagrama de tallo y hojas, diagrama de cajas y bigote, y su histograma y tabla de frecuencias. Analice.
\section{Anàlisis}
Analicelosresultadosdelasecciónanterior,estasdosseccionespueden ir juntas como “Resultados y Análisis”
\section{Conclusiones}
Escriba las conclusiones de sus análisis.

\end{document}
